\section{Methods}
\label{sec:methods}

\subsection{The calculation}

Litres to tonnes to energy to emissions, once per fuel and per gas:
\begin{equation}
E_{f,g} \;=\; Q_f \cdot \rho_f \cdot \mathrm{NCV}_f \cdot \mathrm{EF}_{g},
\label{eq:main}
\end{equation}
where $Q_f$ is the published quantity, $\rho_f$ the density needed to reach a mass basis,
$\mathrm{NCV}_f$ the net calorific value, and $\mathrm{EF}_{g}$ the IPCC rail factor for gas
$g$. There is nothing else in the sector.

The density matters more than it looks. Pakistan Railways reports litres and every IPCC factor
is per unit of energy, so a density is unavoidable, and it is taken from Pakistan State Oil's
own high-speed-diesel material safety data sheet---specific gravity \numrange{0.83}{0.86} at
\SI{16}{\celsius}, midpoint \SI{0.845}{\kilo\gram\per\litre}---rather than from a generic
international default. That spread is the entire uncertainty in the step, and at
\SI{2}{\percent} it is the smallest term in the sector.

\subsection{The furnace-oil rows}

Furnace oil takes its \COtwo{} factor from the fuel's own carbon content in Volume~2 Chapter~2
Table~2.2, because Table~3.4.1 tabulates only diesel and sub-bituminous coal. Its \CHfour{} and
\NtwoO{} factors are the diesel rail values, following §3.4.1.4, which directs that other
traction fuels burnt in the same engines be treated with the same emission properties. Both
choices are recorded on the factor rows rather than inferred from the code.

\subsection{Fossil methane}

Diesel and furnace oil are fossil fuels, so their \CHfour{} takes AR5's fossil value of
\num{30} rather than the biogenic \num{28}---the same treatment as road transport, power and
aviation, and the opposite of waste. A test asserts the \COe{} of every \CHfour{} row against
\num{30} exactly, because the two values differ by seven percent and nothing in the output
distinguishes them.

\subsection{The air-pollutant layer}
\label{sec:air}

Air pollutants take the EMEP/EEA Guidebook's railway chapter \citep{emep2023rail} at
Tier~1: factors per tonne of fuel, on the identical fuel-mass chain as the greenhouse
gases, in kilotonnes of native mass that never enter a \COe{} total. Black carbon is the
chapter's declared fraction (\num{0.65}) of PM\textsubscript{2.5}---a component of the
particulate mass, never an addition to it---and the chapter's one defect is recorded
rather than repaired: its printed 95\% intervals for the three particulate rows do not
bracket their own central values, so those rows carry the printed bounds verbatim and a
conservative $\pm$\SI{100}{\percent}.

\SOtwo{} is stoichiometry, not a factor: twice the fuel's sulfur mass. The chapter says
operator sulfur values are preferable to defaults; Pakistan Railways publishes none, and
Pakistani diesel spans three orders of magnitude---\SI{10}{ppm} Euro-V imports, mandatory
since January 2021, against the \SI{500}{ppm} pre-Euro-V specification and the older
specification's \SI{1}{\percent} ceiling that covers most of the historical series. So
\SOtwo{} ships as an explicit case ladder with the \SI{500}{ppm} specification as the
declared reference, not as one number pretending to precision the fuel supply does not
have.

A Tier~2 diagnostic accompanies the reported Tier~1 layer: fuel split among line-haul
passenger, line-haul freight and shunting (railcars are not occurring---workings are Nil
in every recent year), rescaled to the metered control exactly, as the chapter itself
instructs. Published train-kilometres divide the line-haul fuel; the shunting share is a
declared five-percent assumption, derived from the chapter's own duty rates and
gap-registered, because no shunting activity is published in any form. What the diagnostic
shows is the duty-cycle sensitivity: a nearly-all-line-haul fleet implies Tier~1
understates NO\textsubscript{x} by roughly a fifth and overstates fine particulate mass by
about a sixth, and neither layer is silently swapped for the other.

\subsection{The diesel closure against road transport}
\label{sec:closure}

Pakistan's national high-speed-diesel supply figure covers road, rail, agriculture and
captive generation together. While the road transport sector anchored on that single national
figure, every litre Pakistan Railways burns was at risk of being counted twice---once here and
once as road diesel---and the road sector deducted this module's diesel, litre for litre, by
calling it rather than by holding a copy of the number.

\textbf{That deduction has been retired, and this section records why rather than describing a
mechanism that no longer runs.} Road transport is now anchored on the petroleum industry's own
customer classes, in which Railways is a separate row from Road. The double-count is excluded
by construction at the source, so subtracting the railway account from the road control would
now remove the same fuel twice.

What the railway fuel account does instead is better. Freed from being an internal correction,
it is an \emph{independent cross-source check}: the operator's own metered purchases, from the
Year Book, against the petroleum industry's Railways customer-class total, two organisations
counting the same physical fuel by different routes. The quantity itself is unchanged---
the module's metered diesel is \SI{112.04}{\kilo\tonne}, traction and non-traction
together---and only what the inventory does with it has moved. An internal deduction can only be
arithmetically consistent; an independent comparison can be wrong, which is what makes it
worth reporting. Section~\ref{sec:validation} gives the comparison.

The general point outlives this instance. A closure that becomes structurally impossible has
been superseded by better source data rather than broken, and the right response is to delete
the correction and say so---not to keep a compensating adjustment running against a control
that no longer needs it.

One deduction from road transport's controls does survive, and it is not this sector's:
household generator petrol, owned by the residential module. A test asserts that the railway
account contributes nothing to it---Pakistan Railways burns no petrol, and deducting from a
fuel the operator does not use would be a second error correcting the first.

\subsection{Spatial support and the ODbL route alternative}
\label{sec:spatial}

The following route reconstruction is an ODbL-derived alternative, not the primary package's licensed grid. The primary product retains its compatible administrative support and records unsupported finer placement as unallocated. A source-specific physical rail release must not inherit an unrelated grid merely to eliminate that remainder.

Railway emissions can occur only along a railway, so the allocation is constrained to one.
Line-haul fuel runs along the OpenStreetMap-derived network---\num{7574} active
line-kilometres, \SI{97}{\percent} of the Year Book's \num{7791} route-kilometres---divided
into corridors; shunting fuel sits at the seven divisional yards, weighted by their
published wagon loadings; and each of the three service streams takes its fuel share from
the Tier~2 split, on a surface that reconciles to the national control exactly at national,
provincial and district level on every run.

The honest split of that construction: the network geometry, the yard locations and the
route-kilometre audit are data; the corridor service weights are \emph{declared judgment},
stated basis-by-basis on their data rows, because the operator's public timetable interface
refuses unauthenticated access and no service-frequency statistics are published. The
previous proxy---freight wagons loaded by division, which credited Sindh with fuel burnt
hauling Karachi-origin freight up-country and allocated passenger operations (four times the
freight train-kilometres) by freight geography---is retained as a structural alternative,
and both splits are published. The national total is identical under either; the engine
reads the transcribed share row, and a test pins that row against a fresh spatial
computation so the two cannot drift apart.
